% !TeX root=../main.tex
\chapter{اثبات‌ها و محاسبات تکمیلی گرادیان}
\label{app:ch3-gradient-derivations}

در این پیوست، محاسبات تکمیلی و اثبات روابط
بخش~\ref{cmp:sec:ch3-gradient-calculation} ارائه
می‌شوند. فرض‌ها و نمادگذاری مطابق متن اصلی‌اند،
مگر آنکه خلاف آن تصریح شود.

\section{مشتق تابع زیان نسبت به ناهمانندی‌ها}
\label{app:ch3-gradient-loss-dissimilarities}

در این بخش، ناهمانندی‌های نمونهٔ آموزشی تا مراکز
طبقه‌ها را مطابق
رابطه~\eqref{cmp:eq:ch3-general-center-dissimilarity}
و برای سنجهٔ ثابت $\psi$ با $d_{i,c}$ نشان می‌دهیم.
در مشتق‌گیری این مرحله، این کمیت‌ها مستقل فرض
می‌شوند و وزن طبقه‌ها، برچسب‌ها و پارامترهای مقیاس
توابع زیان ثابت‌اند.

\subsection{مشتق تابع هدف وزن‌دار و حاشیه‌ها}
\label{app:ch3-gradient-loss-margin-chain}

در تابع هدف
رابطه~\eqref{cmp:eq:ch3-training-weighted-loss}،
زیان نمونهٔ $j$ به ناهمانندی‌های همان نمونه وابسته
است؛ ازاین‌رو، برای $j\ne i$،
$\partial\ell_j/\partial d_{i,c}=0$.
در نتیجه
\[
\begin{aligned}
\frac{\partial\mathcal L_{\mathrm{tr}}}{\partial d_{i,c}}
&=\sum_{k=1}^{K}\frac{\omega_k}{|C_k|}
\sum_{j\in C_k}\frac{\partial\ell_j}{\partial d_{i,c}}\\
&=\frac{\omega_{y_i}}{|C_{y_i}|}
\frac{\partial\ell_i}{\partial d_{i,c}},
\end{aligned}
\]
که رابطه~\eqref{cmp:eq:ch3-gradient-loss-distance-coefficient}
را ثابت می‌کند. استقلال در این مرحله صرفاً
قرارداد مشتق جزئی است، نه استقلال ناهمانندی‌ها
نسبت به پارامترهای سیاست. برای شباهت نیز
جایگذاری $d_{i,c}=-s(\mathscr Y_i,\mathcal V_c)$
و قاعدهٔ زنجیره‌ای
رابطه~\eqref{cmp:eq:ch3-gradient-loss-similarity-coefficient}
را نتیجه می‌دهد.

حاشیهٔ تعلق به طبقهٔ $c$ طبق
رابطه~\eqref{cmp:eq:ch3-fixed-representation-margin}
برابر است با
\[
\gamma_{i,c}^{(\psi)}
=\min_{u\ne c}d_{i,u}-d_{i,c}.
\]
اگر نزدیک‌ترین رقیب $r_i(c)$ طبق
رابطه~\eqref{cmp:eq:ch3-gradient-nearest-rival}
یکتا باشد، به دلیل متناهی‌بودن شمار طبقه‌ها و
شکاف مثبت با دیگر رقبا، شاخهٔ کمینه در یک
همسایگی از بردار ناهمانندی‌ها ثابت می‌ماند.
برای هر $r=1,\ldots,K$ داریم:
\begin{equation}\label{eq:app-ch3-gradient-margin-distance}
\frac{\partial\gamma_{i,c}^{(\psi)}}{\partial d_{i,r}}
=\mathbf1\{r=r_i(c)\}-\mathbf1\{r=c\}.
\end{equation}
همچنین از تعریف حاشیهٔ دوتایی در
رابطه~\eqref{cmp:eq:ch3-training-pairwise-margin}،
برای $u\ne y_i$ نتیجه می‌شود:
\begin{equation}\label{eq:app-ch3-gradient-pairwise-distance}
\frac{\partial m_{i,u}^{(\psi)}}{\partial d_{i,r}}
=\mathbf1\{r=u\}-\mathbf1\{r=y_i\}.
\end{equation}
رابطهٔ اخیر به یکتایی هیچ رقیبی نیاز ندارد.

\subsection{استخراج مشتق سه تابع زیان}
\label{app:ch3-gradient-loss-three-derivatives}

\paragraph{زیان لگاریتمی.}
در حالت یکتایی نزدیک‌ترین رقیب برای همهٔ
طبقه‌ها، قاعدهٔ زنجیره‌ای و
روابط~\eqref{cmp:eq:ch3-training-margin-nll-derivative}
و~\eqref{eq:app-ch3-gradient-margin-distance}
نتیجه می‌دهند:
\[
\begin{aligned}
\frac{\partial\ell_i}{\partial d_{i,r}}
&=\sum_{c=1}^{K}
\frac{\partial\ell_i}{\partial\gamma_{i,c}^{(\psi)}}
\frac{\partial\gamma_{i,c}^{(\psi)}}{\partial d_{i,r}}\\
&=\sum_{c=1}^{K}
\frac{\widehat p_{i,c}^{(\psi,\tau_\psi)}-Y_{ic}}{\tau_\psi}
\bigl(\mathbf1\{r=r_i(c)\}-\mathbf1\{r=c\}\bigr).
\end{aligned}
\]
بنابراین رابطه~\eqref{cmp:eq:ch3-gradient-nll-distance}
برقرار است.

\paragraph{زیان سیگموئیدی.}
در حالت یکتایی نزدیک‌ترین رقیب طبقهٔ واقعی، از
روابط~\eqref{cmp:eq:ch3-training-sigmoid-derivative}
و~\eqref{eq:app-ch3-gradient-margin-distance}
به دست می‌آید:
\[
\begin{aligned}
\frac{\partial\ell_i}{\partial d_{i,r}}
&=\frac{\partial\ell_i}{\partial\gamma_{i,y_i}^{(\psi)}}
\frac{\partial\gamma_{i,y_i}^{(\psi)}}{\partial d_{i,r}}\\
&=-\frac{\ell_i(1-\ell_i)}{\tau}
\bigl(\mathbf1\{r=r_i(y_i)\}-\mathbf1\{r=y_i\}\bigr).
\end{aligned}
\]
این همان رابطه~\eqref{cmp:eq:ch3-gradient-sigmoid-distance}
است.

\paragraph{زیان سافت‌پلاس.}
از قاعدهٔ زنجیره‌ای روی همهٔ حاشیه‌های دوتایی
و رابطه~\eqref{eq:app-ch3-gradient-pairwise-distance}
داریم:
\[
\frac{\partial\ell_i}{\partial d_{i,r}}
=\sum_{u\ne y_i}
\frac{\partial\ell_i}{\partial m_{i,u}^{(\psi)}}
\bigl(\mathbf1\{r=u\}-\mathbf1\{r=y_i\}\bigr).
\]
برای $r\ne y_i$ فقط جملهٔ $u=r$ باقی
می‌ماند و برای $r=y_i$ همهٔ جمله‌ها با علامت
منفی جمع می‌شوند؛ بنابراین
رابطه~\eqref{cmp:eq:ch3-gradient-softplus-distance}
حاصل می‌شود. مشتق هر جمله نسبت به حاشیه
در رابطه~\eqref{cmp:eq:ch3-training-softplus-derivative}
محاسبه شده است و به انتخاب نزدیک‌ترین رقیب
وابسته نیست.

\subsection{تساوی نزدیک‌ترین رقبا و مشتق جهت‌دار}
\label{app:ch3-gradient-loss-ties}

اگر نزدیک‌ترین رقیب یکتا نباشد، مجموعهٔ اندیس
رقبای فعال را در این زیربخش با
\[
\mathcal R_i(c)=
\{u\ne c:\ d_{i,u}=\min_{r\ne c}d_{i,r}\}
\]
نشان می‌دهیم. برای جهت تغییر
$\boldsymbol v=(v_1,\ldots,v_K)$ در فضای
ناهمانندی‌ها، مشتق جهت‌دار یک‌طرفهٔ حاشیه
برابر است با
\begin{equation}\label{eq:app-ch3-gradient-margin-directional}
\lim_{\varepsilon\downarrow0}
\frac{\min_{u\ne c}(d_{i,u}+\varepsilon v_u)
-(d_{i,c}+\varepsilon v_c)-\gamma_{i,c}^{(\psi)}}
{\varepsilon}
=\min_{u\in\mathcal R_i(c)}v_u-v_c.
\end{equation}
زیرا برای $\varepsilon$ مثبت و به‌اندازهٔ کافی
کوچک، هیچ رقیب غیرفعالی به دلیل شکاف مثبت
نمی‌تواند به کمینه برسد. اگر همهٔ رقبا فعال
باشند، رابطه مستقیماً از تعریف کمینه به دست
می‌آید.

در حالت تساوی چند رقیب، عبارت سمت راست
به‌طور کلی تابعی خطی از جهت نیست و حاشیه
مشتق معمول ندارد. مشتق جهت‌دار زیان لگاریتمی
از جمع مشتق‌های جهت‌دار حاشیه‌ها با ضرایب
رابطه~\eqref{cmp:eq:ch3-training-margin-nll-derivative}
و مشتق جهت‌دار زیان سیگموئیدی از ضرب مشتق
جهت‌دار حاشیهٔ طبقهٔ واقعی در
رابطه~\eqref{cmp:eq:ch3-training-sigmoid-derivative}
حاصل می‌شود. ممکن است ترکیب این مشتق‌ها در
موارد خاص خطی شود؛ بنابراین عدم مشتق‌پذیری
حاشیه لزوماً به عدم مشتق‌پذیری زیان نمی‌انجامد.

برای محاسبهٔ عددی می‌توان یک عضو فعال را
با قراردادی ثابت، مانند کوچک‌ترین اندیس فعال،
انتخاب کرد و در روابط متن به جای رقیب یکتا
قرار داد. همچنین در سطح هر حاشیه، می‌توان
بردارهای مشتق شاخه‌های فعال را با وزن‌های
نامنفی و مجموع یک ترکیب کرد؛ در این صورت
$\mathbf1\{r=r_i(c)\}$ با وزن مربوط به
رقیب $r$ جایگزین می‌شود (وزن رقبا خارج از
مجموعهٔ فعال صفر است).
این‌ها قراردادهای انتشار مشتق در نقطهٔ تساوی‌اند؛
نباید آن‌ها را با مشتق معمول یا مشتق جهت‌دار
در تمام جهت‌ها یکی دانست. همچنین در زیان
لگاریتمی، انتخاب‌های مستقل شاخه‌ها برای حاشیه‌های
مختلف لزوماً از یک همسایگی مشتق‌پذیر مشترک
ناشی نمی‌شوند.

\subsection{ناوردایی انتقال مشترک و خاصیت مجموع صفر}
\label{app:ch3-gradient-loss-zero-sum}

اگر به تمام ناهمانندی‌های نمونهٔ $i$ مقدار
ثابت $t$ افزوده شود، تمام حاشیه‌های دوتایی
و حاشیه‌های تعلق بدون تغییر می‌مانند، زیرا
اختلاف ناهمانندی‌ها ثابت می‌ماند. در نتیجه
بردار احتمال تخمینی برچسب و هر سه زیان نیز
ثابت می‌مانند. در نقاط مشتق‌پذیر، مشتق این
همانی نسبت به $t$ در $t=0$ برابر است با
\[
0=\left.\frac{\mathrm d}{\mathrm dt}
\ell_i(d_{i,1}+t,\ldots,d_{i,K}+t)\right|_{t=0}
=\sum_{r=1}^{K}\frac{\partial\ell_i}{\partial d_{i,r}}.
\]
در قرارداد انتخاب شاخهٔ فعال نیز مجموع
مؤلفه‌های بردار
$\mathbf1\{r=r_i(c)\}-\mathbf1\{r=c\}$
صفر است؛ همین ویژگی برای ترکیب وزن‌دار
شاخه‌ها برقرار می‌ماند. برای سافت‌پلاس، صفر
بودن مجموع مستقیماً از
رابطه~\eqref{cmp:eq:ch3-gradient-softplus-distance}
نتیجه می‌شود. در نهایت با ضرب در وزن ثابت
$\omega_{y_i}/|C_{y_i}|$،
رابطه~\eqref{cmp:eq:ch3-gradient-loss-zero-sum}
به دست می‌آید.

\section{مشتق سنجه‌های ناهمانندی و شباهت}
\label{app:ch3-gradient-measures}

در این بخش، روابط مشتق‌گیری
زیربخش~\ref{cmp:subsec:ch3-gradient-measures}
را اثبات می‌کنیم. تعریف سنجه‌ها همان تعریف‌های
زیربخش~\ref{cmp:subsec:ch3-decision-dissimilarities}
است و تمام لگاریتم‌ها طبیعی‌اند.

\subsection{مشتق‌گیری روی سیمپلکس احتمال}
\label{app:ch3-gradient-measures-simplex}

برای توزیع‌های دارای مؤلفه‌های اکیداً مثبت،
فرمول‌های تعریف سنجه را به همسایگی مثبت
در فضای $\mathbb R^K\times\mathbb R^K$ امتداد
می‌دهیم و نسبت به مختصات مستقل مشتق می‌گیریم.
اگر $P$ و $Q$ بر حسب پارامترهای سیاست
تغییر کنند و مجموع مؤلفه‌های هر یک همواره
یک باشد، از مشتق این دو قید داریم
$\sum_u\mathrm dP(u)=\sum_u\mathrm dQ(u)=0$.
برای سنجهٔ مشتق‌پذیر $d_\psi$ در این همسایگی،
قاعدهٔ زنجیره‌ای نتیجه می‌دهد:
\[
\mathrm d d_\psi(P,Q)
=(\nabla_P d_\psi(P,Q))^{\mathsf T}\mathrm dP
+(\nabla_Q d_\psi(P,Q))^{\mathsf T}\mathrm dQ.
\]
افزودن هر مضرب ثابت از بردار همه‌یک به
هر یک از دو گرادیان، مقدار این دیفرانسیل را
تغییر نمی‌دهد، زیرا تغییرات توزیع‌ها مجموع
صفر دارند. ازاین‌رو مشتق‌های فضای محیطی برای
انتشار مشتق در مسیرهای مجاز سیمپلکس کافی‌اند؛
محاسبهٔ گرادیان مقید جداگانه لازم نیست.

\subsection{واگرایی کولبک--لایبلر}
\label{app:ch3-gradient-measures-kl}

از تعریف~\eqref{cmp:eq:ch3-KL-definition} داریم:
\[
d_{\mathrm{KL}}(P,Q)
=\sum_{u=1}^{K}\bigl(P(u)\log P(u)-P(u)\log Q(u)\bigr).
\]
برای $P(u),Q(u)>0$، تنها جملهٔ متناظر با
مؤلفهٔ $u$ در مشتق جزئی آن مؤلفه دخالت دارد:
\[
\begin{aligned}
\frac{\partial d_{\mathrm{KL}}}{\partial P(u)}
&=\log P(u)+1-\log Q(u),\\
\frac{\partial d_{\mathrm{KL}}}{\partial Q(u)}
&=-P(u)/Q(u).
\end{aligned}
\]
با ترکیب لگاریتم‌ها رابطهٔ
\eqref{cmp:eq:ch3-gradient-kl-partials} حاصل می‌شود.

\subsection{واگرایی جنسن--شنون}
\label{app:ch3-gradient-measures-js}

با استفاده از تعریف
رابطه~\eqref{cmp:eq:ch3-decision-js} و
هویت آنتروپی در رابطه~\eqref{cmp:eq:ch3-JSD-entropy}
می‌نویسیم:
\[
d_{\mathrm{JS}}(P,Q)
=H\!\left(\frac{P+Q}{2}\right)-\frac{H(P)+H(Q)}2 .
\]
در این رابطه $H(\cdot)$ تابع آنتروپی است.
برای مؤلفه‌های مثبت داریم:
\[
\begin{aligned}
\frac{\partial d_{\mathrm{JS}}}{\partial P(u)}
&=-\frac12\left(\log\frac{P(u)+Q(u)}2+1\right)
+\frac12(\log P(u)+1)\\
&=\frac12\log\frac{2P(u)}{P(u)+Q(u)}.
\end{aligned}
\]
با جابه‌جایی $P$ و $Q$ و تقارن JSD، مشتق
نسبت به آرگومان دوم نیز حاصل می‌شود. در نتیجه
رابطه~\eqref{cmp:eq:ch3-gradient-js-partials}
اثبات می‌شود.

\subsection{متریک جذر جنسن--شنون و شباهت}
\label{app:ch3-gradient-measures-sqrt}

در نقاط مشتق‌پذیری JSD که مقدار آن مثبت
است، مشتق تابع جذر برابر
$1/(2\sqrt{d_{\mathrm{JS}}(P,Q)})$ است.
با قاعدهٔ زنجیره‌ای و تعریف
رابطه~\eqref{cmp:eq:ch3-decision-js-metric}،
رابطه~\eqref{cmp:eq:ch3-gradient-sqrt-js-partials}
حاصل می‌شود. همچنین تعریف شباهت در
رابطه~\eqref{cmp:eq:ch3-decision-js-similarity}
نگاشتی آفین از JSD است؛ بنابراین مشتق
آن نسبت به هر آرگومان، مشتق JSD ضرب در
$-1/\log2$ است و
رابطه~\eqref{cmp:eq:ch3-gradient-js-similarity-partials}
برقرار می‌شود.

\subsection{شرایط مرزی و تکیه‌گاه ثابت}
\label{app:ch3-gradient-measures-boundary}

قرارداد $0\log0=0$ برای مقدار آنتروپی
به معنای وجود مشتق متناهی در صفر نیست.
فرض کنیم تکیه‌گاه هر یک از دو توزیع در
همسایگی پارامترها ثابت باشد؛ مؤلفه‌های صفر
در این همسایگی ثابت‌اند و از محاسبهٔ مشتق
حذف می‌شوند. نباید حاصل‌ضرب صوری
$0\cdot\log0$ را در فرمول مشتق به کار برد.

\paragraph{واگرایی کولبک--لایبلر.}
اگر $P(u)>0$ و $Q(u)=0$ باشد، واگرایی
نامتناهی است و مشتق متناهی تابع هدف مبتنی
بر آن قابل محاسبه نیست. بنابراین برای KL
متناهی، تکیه‌گاه $P$ باید زیرمجموعهٔ
تکیه‌گاه $Q$ باشد. در تکیه‌گاه ثابت:
اگر هر دو مؤلفه مثبت باشند از
رابطه~\eqref{cmp:eq:ch3-gradient-kl-partials}
استفاده می‌شود؛ اگر $P(u)=0$ و $Q(u)>0$
باشد، جملهٔ متناظر همواره صفر و مشتق آن
نسبت به مؤلفهٔ مثبت $Q(u)$ نیز صفر است.
اگر هر دو مؤلفه صفر ساختاری باشند، سهم
آن‌ها در دیفرانسیل صفر است.

\paragraph{واگرایی جنسن--شنون و شباهت.}
در JSD تکیه‌گاه‌ها می‌توانند متفاوت باشند.
برای $P(u)>0$ و $Q(u)=0$ ساختاری،
سهم مؤلفهٔ $u$ برابر $P(u)\log2/2$ و
مشتق نسبت به $P(u)$ برابر $\log2/2$ است.
حالت $P(u)=0$ ساختاری و $Q(u)>0$ متقارن
است. اگر هر دو مؤلفه صفر ساختاری باشند،
سهم و تغییر آن صفر است. در هر دو سنجهٔ
JSD و KL، هنگامی که مؤلفهٔ صفر آرگومان
نخست در برابر مؤلفهٔ مثبت آرگومان دوم
از صفر افزایش می‌یابد، عبارت $x\log x$
دارای مشتق راست نامتناهی منفی است.
برای JSD همین وضعیت با تعویض آرگومان‌ها
نیز رخ می‌دهد. اگر هر دو مؤلفه صفر باشند
و بتوانند هم‌زمان تغییر کنند، رفتار به مسیر
تغییر وابسته است: برای JSD سهم مؤلفه در
مسیر $(P(u),Q(u))=(ta,tb)$، با $t>0$،
دقیقاً $t$ برابر سهم متناظر با $(a,b)$ است؛
این سهم نسبت به $(a,b)$ عموماً خطی نیست.
بنابراین متناهی بودن JSD وجود مشتق معمول
در همهٔ جهت‌های مرزی را تضمین نمی‌کند.
نتایج شباهت از رابطهٔ آفین آن با JSD حاصل
می‌شوند.

\paragraph{رفتار متریک جذر در برابری دو توزیع.}
برای نشان دادن مسئله حتی در داخل سیمپلکس،
$P=Q$ را با مؤلفه‌های مثبت در نظر بگیرید
و بردار غیرصفر $v$ را با $\sum_u v(u)=0$
انتخاب کنید. برای $t$ با قدر مطلق کافی کوچک،
$P+tv$ یک توزیع احتمال است. از بسط تیلور
آنتروپی، با مشتق دوم $-1/P(u)$ در هر
مؤلفه، به دست می‌آید:
\[
d_{\mathrm{JS}}(P+tv,P)
=\frac{t^2}{8}\sum_{u=1}^{K}\frac{v(u)^2}{P(u)}
+o(t^2).
\]
در نتیجه
\[
d_{\sqrt{\mathrm{JS}}}(P+tv,P)
=|t|\sqrt{\frac18\sum_{u=1}^{K}
\frac{v(u)^2}{P(u)}}+o(|t|).
\]
مشتق‌های چپ و راست نسبت به $t$ متفاوت‌اند؛
پس متریک جذر JSD در نقطهٔ برابری دو
توزیع داخلی مشتق معمول ندارد. در یک مسیر
خاص پارامترهای سیاست ممکن است دو توزیع
با هم ثابت بمانند یا تغییر مرتبهٔ نخست
آن‌ها یکسان باشد؛ ازاین‌رو عدم مشتق‌پذیری
سنجه به‌تنهایی تعیین‌کنندهٔ مشتق تابع مرکب
در تمام مسیرها نیست.

\subsection{هموارسازی و انتشار مشتق}
\label{app:ch3-gradient-measures-smoothing}

هموارسازی مشترک
رابطه~\eqref{cmp:eq:ch3-stable-smoothed-representations}
برای توزیع‌های عمومی $P,Q$ با بردار یکنواخت
$U$ و پارامتر ثابت $0<\eta<1$، مؤلفه‌های
$P^{(\eta)}=(1-\eta)P+\eta U$ و
$Q^{(\eta)}=(1-\eta)Q+\eta U$ را در
بازهٔ $[\eta/K,1]$ قرار می‌دهد.
در نتیجه:
\[
\begin{aligned}
0\le D_{\mathrm{KL}}(P^{(\eta)}\|Q^{(\eta)})
&\le-\sum_u P^{(\eta)}(u)\log Q^{(\eta)}(u)\\
&\le\log(K/\eta)<\infty.
\end{aligned}
\]
کران پایین از نامنفی بودن KL و کران بالا
از $\sum_u P^{(\eta)}(u)\log P^{(\eta)}(u)\le0$
و حد پایین مؤلفه‌های $Q^{(\eta)}$ به دست می‌آید.

با ثابت بودن $\eta$ و $U$ داریم:
\[
\frac{\partial P^{(\eta)}(u)}{\partial P(v)}
=(1-\eta)\mathbf1\{u=v\},
\qquad
\frac{\partial Q^{(\eta)}(u)}{\partial Q(v)}
=(1-\eta)\mathbf1\{u=v\}.
\]
ازاین‌رو برای تابع سنجهٔ مشتق‌پذیر در
ورودی‌های هموارشده:
\[
\begin{aligned}
\nabla_P[d_\psi(P^{(\eta)},Q^{(\eta)})]
&=(1-\eta)\nabla_{P^{(\eta)}}d_\psi(P^{(\eta)},Q^{(\eta)}),\\
\nabla_Q[d_\psi(P^{(\eta)},Q^{(\eta)})]
&=(1-\eta)\nabla_{Q^{(\eta)}}d_\psi(P^{(\eta)},Q^{(\eta)}).
\end{aligned}
\]
رابطه برای شباهت نیز به همین صورت برقرار
است. از آفین بودن نگاشت هموارسازی نتیجه
می‌شود که میانگین بازنمایی‌های هموارشدهٔ
اعضای هر طبقه همان مرکز هموارشدهٔ آن طبقه
است. این سازگاری اجازه می‌دهد در زنجیرهٔ
مشتق از همان ساختار میانگین مراکز استفاده
شود، اما عامل $(1-\eta)$ باید دقیقاً یک‌بار
در هر مسیر وابستگی لحاظ شود. هموارسازی
تابع هدف را تغییر می‌دهد و در مورد متریک
جذر، چون
$P^{(\eta)}=Q^{(\eta)}$ اگر و تنها اگر $P=Q$،
مسئلهٔ مشتق در برابری دو توزیع باقی می‌ماند.

\section{انتقال مشتق به بازنمایی‌ها و مراکز طبقه‌ها}
\label{app:ch3-gradient-centers-representations}

در این بخش روابط
زیربخش~\ref{cmp:subsec:ch3-gradient-centers-representations}
را از تعریف مراکز، حساسیت‌های تابع هدف و
قاعدهٔ زنجیره‌ای استخراج می‌کنیم. داده، برچسب‌ها،
مجموعهٔ اعضای هر طبقه و وزن طبقه‌ها ثابت‌اند
و شرایط مشتق‌پذیری سنجه‌ها مطابق
بخش~\ref{app:ch3-gradient-measures} است.

\subsection{مشتق نگاشت تشکیل مراکز}
\label{app:ch3-gradient-center-map}

طبق رابطه~\eqref{cmp:eq:ch3-class-center}،
مرکز طبقهٔ $c$ میانگین حسابی بازنمایی‌های
اعضای آن است. برای دو اندیس مؤلفهٔ
$u,v\in\{1,\ldots,K\}$ داریم:
\[
\frac{\partial\mathcal V_c(u)}
{\partial\mathscr Y_s(v)}
=\frac{\mathbf1\{y_s=c\}}{|C_c|}
\mathbf1\{u=v\}.
\]
بنابراین ژاکوبین نگاشت مرکز نسبت به
مختصات بازنمایی نمونه برابر است با:
\begin{equation}\label{eq:app-ch3-gradient-center-jacobian}
\boxed{
\frac{\partial\mathcal V_c}{\partial\mathscr Y_s}
=\frac{\mathbf1\{y_s=c\}}{|C_c|}I_K.
}
\end{equation}
در اینجا $I_K$ ماتریس همانی مرتبهٔ $K$ است.
رابطه در فضای مختصات محیطی نوشته شده است
و روی تغییرات مجاز سیمپلکس نیز معتبر است.
با مشتق‌گیری از میانگین متناهی نسبت به
پارامترهای سیاست، برای هر مؤلفهٔ مرکز داریم:
\begin{equation}\label{cmp:eq:ch3-margin-center-policy-derivative}
\boxed{
\nabla_\theta\mathcal V_c(u)
=\frac1{|C_c|}\sum_{j\in C_c}
\nabla_\theta\mathscr Y_j(u).
}
\end{equation}
به‌طور معادل، ژاکوبین مرکز نسبت به پارامترها
میانگین ژاکوبین‌های بازنمایی اعضای آن است.
این نتیجه به نوع سنجه یا تابع زیان وابسته نیست.

\subsection{مشتق جزئی تابع هدف نسبت به مراکز}
\label{app:ch3-gradient-center-accumulation}

ابتدا تابع هدف آموزش را به‌صورت تابعی از
بازنمایی‌ها و مراکز مستقل در نظر می‌گیریم.
در این قرارداد، تغییر مرکز $\mathcal V_c$
تنها ناهمانندی‌های $d_{i,c}^{(\psi)}$ را برای
$i=1,\ldots,n$ تغییر می‌دهد؛ تمام بازنمایی‌ها
و سایر مراکز ثابت‌اند. با استفاده از تعریف
حساسیت در
رابطه~\eqref{cmp:eq:ch3-gradient-loss-distance-coefficient}
و قاعدهٔ زنجیره‌ای داریم:
\[
\begin{aligned}
\nabla_{\mathcal V_c}\mathcal L_{\mathrm{tr}}
&=\sum_{i=1}^{n}
\frac{\partial\mathcal L_{\mathrm{tr}}}
{\partial d_{i,c}^{(\psi)}}
\nabla_Q d_\psi(\mathscr Y_i,\mathcal V_c)\\
&=\sum_{i=1}^{n}g_{i,c}^{(\psi)}
\nabla_Q d_\psi(\mathscr Y_i,\mathcal V_c).
\end{aligned}
\]
این همان رابطه~\eqref{cmp:eq:ch3-gradient-center-accumulation}
است. وزن هر نمونه قبلاً در $g_{i,c}^{(\psi)}$
لحاظ شده است و نباید هنگام تجمیع دوباره
اعمال شود.

\subsection{گرادیان کل نسبت به بازنمایی}
\label{app:ch3-gradient-representation-total}

با ثابت فرض کردن مراکز، تنها ناهمانندی‌های
نمونهٔ $s$ نسبت به تمام مراکز به بازنمایی
$\mathscr Y_s$ وابسته‌اند. بنابراین سهم مستقیم
گرادیان برابر است با
\[
\sum_{c=1}^{K}g_{s,c}^{(\psi)}
\nabla_P d_\psi(\mathscr Y_s,\mathcal V_c).
\]
اکنون مراکز را تابع بازنمایی‌ها قرار می‌دهیم.
مطابق قاعدهٔ زنجیره‌ای و
رابطه~\eqref{eq:app-ch3-gradient-center-jacobian}،
سهم غیرمستقیم برابر است با
\[
\begin{aligned}
\sum_{c=1}^{K}
\left(\frac{\partial\mathcal V_c}
{\partial\mathscr Y_s}\right)^{\mathsf T}
\nabla_{\mathcal V_c}\mathcal L_{\mathrm{tr}}
&=\sum_{c=1}^{K}
\frac{\mathbf1\{y_s=c\}}{|C_c|}
\nabla_{\mathcal V_c}\mathcal L_{\mathrm{tr}}\\
&=\frac1{|C_{y_s}|}
\nabla_{\mathcal V_{y_s}}\mathcal L_{\mathrm{tr}}.
\end{aligned}
\]
جمع این دو سهم،
رابطه~\eqref{cmp:eq:ch3-gradient-representation-total}
را نتیجه می‌دهد. در این رابطه سهم مرکز شامل
زیان خود نمونه نیز هست؛ دو سهم، دو مسیر
متفاوت وابستگی همان تابع هدف‌اند، نه دو
تابع هدف جداگانه.

از آنجا که بازنمایی در محاسبات پیش‌رونده
گشت به‌صورت بردار سطری نوشته می‌شود، در این
قرارداد گرادیان $\mathbf g_s$ ستونی است؛
در نتیجه دیفرانسیل تابع هدف را می‌توان به
صورت زیر نوشت:
\[
\mathrm d\mathcal L_{\mathrm{tr}}
=\sum_{s=1}^{n}(\mathrm d\mathscr Y_s)\mathbf g_s.
\]
برای شباهت، همین استخراج با جایگزینی
$d_\psi$ و $g_{i,c}^{(\psi)}$ به‌ترتیب با
$s$ و $g_{i,c}^{(s)}$ برقرار است.

\subsection{سازگاری با هموارسازی}
\label{app:ch3-gradient-centers-smoothing}

اگر از نگاشت هموارسازی
رابطه~\eqref{cmp:eq:ch3-stable-smoothed-representations}
استفاده کنیم، میانگین بازنمایی‌های هموارشده
دقیقاً مرکز هموارشده است. مطابق
بخش~\ref{app:ch3-gradient-measures-smoothing}،
دیفرانسیل هر ورودی هموارشده برابر
$(1-\eta)$ برابر دیفرانسیل ورودی خام است.
می‌توان ابتدا گرادیان را نسبت به بازنمایی‌های
هموارشده، شامل هر دو مسیر وابستگی، محاسبه کرد
و سپس حاصل را در $(1-\eta)$ ضرب کرد:
\[
\begin{aligned}
\mathbf g_s
=(1-\eta)\Bigg[
&\sum_{c=1}^{K}g_{s,c}^{(\psi)}
\nabla_P d_\psi(\mathscr Y_s^{(\eta)},\mathcal V_c^{(\eta)})\\
&+\frac1{|C_{y_s}|}\sum_{i=1}^{n}
g_{i,y_s}^{(\psi)}
\nabla_Q d_\psi(\mathscr Y_i^{(\eta)},\mathcal V_{y_s}^{(\eta)})
\Bigg].
\end{aligned}
\]
در این فرمول حساسیت‌ها نیز در ناهمانندی‌های
هموارشده ارزیابی می‌شوند و $\mathbf g_s$
گرادیان نسبت به بازنمایی خام است. به‌طور
معادل، می‌توان سنجهٔ مرکب با نگاشت هموارسازی
را مستقیماً نسبت به ورودی خام مشتق گرفت؛
در آن صورت عامل $(1-\eta)$ در خود مشتق
سنجه وجود دارد و نباید دوباره ضرب شود.
این استدلال برای شباهت نیز برقرار است.

\subsection{صورت خاص مشتق امتیاز مبتنی بر جنسن--شنون}
\label{app:ch3-gradient-centers-js-special}

برای تطبیق با تحلیل‌های بعدی زیان سیگموئیدی،
$P_i=\mathscr Y_i^{(\eta)}$،
$V_c=\mathcal V_c^{(\eta)}$ و
$s_{i,c}=s(P_i,V_c)$ را در نظر می‌گیریم.
در ناحیهٔ مثبت، از
رابطه~\eqref{cmp:eq:ch3-gradient-js-similarity-partials}
داریم:
\begin{equation}\label{cmp:eq:ch3-margin-JS-partials}
\begin{aligned}
\frac{\partial s(P_i,V_c)}{\partial P_i(u)}
&=-\frac1{2\log2}\log\frac{2P_i(u)}{P_i(u)+V_c(u)},\\
\frac{\partial s(P_i,V_c)}{\partial V_c(u)}
&=-\frac1{2\log2}\log\frac{2V_c(u)}{P_i(u)+V_c(u)}.
\end{aligned}
\end{equation}
با قاعدهٔ زنجیره‌ای به دست می‌آید:
\begin{equation}\label{cmp:eq:ch3-margin-score-policy-gradient}
\begin{aligned}
\nabla_\theta s_{i,c}
=-\frac1{2\log2}\sum_{u=1}^{K}\Bigg[
&\log\frac{2P_i(u)}{P_i(u)+V_c(u)}\nabla_\theta P_i(u)\\
+&\log\frac{2V_c(u)}{P_i(u)+V_c(u)}\nabla_\theta V_c(u)
\Bigg].
\end{aligned}
\end{equation}
در اینجا مشتق‌های ورودی‌های هموارشده طبق
بخش~\ref{app:ch3-gradient-centers-smoothing}
محاسبه می‌شوند.

برای صورت سافت‌پلاس مبتنی بر شباهت جنسن--شنون
در تحلیل‌های بعدی، با نمادگذاری همان تحلیل
$z=\operatorname{vec}(Z)$،
$U_i=\mathscr Y_i$،
$V_r=\mathcal V_r$ و $S_{i,r}=s(U_i,V_r)$،
از مشتق امتیاز بالا و
رابطه~\eqref{cmp:eq:ch3-margin-center-policy-derivative}
نتیجه می‌شود:
\begin{equation}\label{cmp:eq:ch3-softplus-js-score-chain}
\begin{aligned}
\nabla_z S_{i,r}
=-\frac1{2\log2}\sum_{u=1}^{K}\Bigg[
&\log\frac{2U_i(u)}{U_i(u)+V_r(u)}\,\nabla_z U_i(u)\\
+&\log\frac{2V_r(u)}{U_i(u)+V_r(u)}
\,\frac1{|C_r|}\sum_{j\in C_r}\nabla_z U_j(u)
\Bigg].
\end{aligned}
\end{equation}
این صورت ویژهٔ همان دو مسیر وابستگی است و
برای مؤلفه‌های مثبت نوشته شده است؛ در
تکیه‌گاه ثابت، حذف مؤلفه‌های ساختاری صفر
مطابق بخش~\ref{app:ch3-gradient-measures-boundary}
انجام می‌شود. در صورت هموارسازی باید عامل
مشتق نگاشت هموارسازی مطابق بخش پیشین
لحاظ شود. این رابطه در هر دو خانوادهٔ سیاست
مشترک و شرطی به طبقه معتبر است.

\section{مشتق‌گیری از گشت دقیق}
\label{app:ch3-gradient-walk}

در این بخش روابط
زیربخش~\ref{cmp:subsec:ch3-gradient-exact-walk}
را اثبات می‌کنیم. داده، برچسب‌ها، فضای حالت،
ماسک‌ها، وزن‌های عمودی و پارامترهای
$\beta,\lambda,H$ ثابت‌اند. در نتیجه احتمال‌های
عمودی، ضرایب اصلاح وابسته به مبدأ،
ضرایب $\alpha_t^{(H)}$ و ماتریس تجمیع برچسب
به پارامترهای قابل‌آموزش سیاست وابسته نیستند.
مشتق‌ها در نقاطی که زنجیرهٔ تابع هدف مشتق‌پذیر
است محاسبه می‌شوند؛ در نقاط تساوی رقبا با
قرارداد مشتق‌گیری بخش‌های پیشین، همان روابط
انتشار حساسیت منتخب را انجام می‌دهند.

\subsection{دیفرانسیل توان ماتریس و بازنمایی سقف‌دار}
\label{app:ch3-gradient-walk-differential}

برای ماتریس مربعی $P_{\mathcal S_s}$،
از قاعدهٔ ضرب ماتریس‌ها و با قرارداد
$P_{\mathcal S_s}^0=I$ داریم:
\begin{equation}\label{eq:app-ch3-power-differential}
\mathrm d(P_{\mathcal S_s}^q)
=\sum_{j=0}^{q-1}P_{\mathcal S_s}^j
(\mathrm dP_{\mathcal S_s})P_{\mathcal S_s}^{q-1-j},
\qquad q\ge1.
\end{equation}
برای $q=1$ رابطه بدیهی است. اگر برای $q$
برقرار باشد، مشتق حاصل‌ضرب
$P_{\mathcal S_s}^qP_{\mathcal S_s}$،
با ضرب جمله‌های جمع در $P_{\mathcal S_s}$
و افزودن $P_{\mathcal S_s}^q\mathrm dP_{\mathcal S_s}$،
همین رابطه را برای $q+1$ نتیجه می‌دهد.
برای $q=0$ مشتق صفر است.
با جایگذاری
$\nu_{s,t}=\rho_sP_{\mathcal S_s}^{t-1}$
در رابطه~\eqref{cmp:eq:ch3-matrix-forward-recursion}،
و با استفاده از مشتق میانگین مراکز در
بخش~\ref{app:ch3-gradient-center-map}، به دست می‌آید:
\begin{equation}\label{eq:app-ch3-center-differential}
\begin{aligned}
\mathrm d\mathscr Y_s
={}&(\mathrm d\rho_s)
\sum_{t=1}^{H}\alpha_t^{(H)}
P_{\mathcal S_s}^{t-1}\mathsf M_s^{(Y)}\\
&+\rho_s\sum_{t=2}^{H}\alpha_t^{(H)}
\sum_{j=0}^{t-2}P_{\mathcal S_s}^j
(\mathrm dP_{\mathcal S_s})
P_{\mathcal S_s}^{t-2-j}\mathsf M_s^{(Y)},\\
\mathrm d\mathcal V_c
={}&|C_c|^{-1}\sum_{s\in C_c}\mathrm d\mathscr Y_s.
\end{aligned}
\end{equation}
این صورت علاوه بر انتشار گرادیان برای محاسبهٔ
ژاکوبین بازنمایی‌ها و مراکز نیز کاربرد دارد.
اگر $\beta$ یا $\lambda$ به‌عنوان پارامتر
قابل‌آموزش انتخاب شوند، باید جمله‌های حاصل
از دیفرانسیل $\alpha_t^{(H)}$ نیز افزوده شوند؛
این حالت جزو قرارداد فعلی آموزش نیست.

\subsection{استخراج بازگشت پس‌روندهٔ حساسیت}
\label{app:ch3-gradient-walk-adjoint}

از بخش~\ref{app:ch3-gradient-representation-total}
می‌دانیم دیفرانسیل تابع هدف برابر مجموع
$(\mathrm d\mathscr Y_s)\mathbf g_s$
روی همهٔ مبدأهاست. برای مبدأ ثابت $s$،
دیفرانسیل بازنمایی پیش‌رونده برابر
\[
\mathrm d\mathscr Y_s
=\sum_{t=1}^{H}\alpha_t^{(H)}
(\mathrm d\nu_{s,t})\mathsf M_s^{(Y)}
\]
است و دیفرانسیل انتقال از رابطهٔ
\[
\mathrm d\nu_{s,t+1}
=(\mathrm d\nu_{s,t})P_{\mathcal S_s}
+\nu_{s,t}\mathrm dP_{\mathcal S_s},
\qquad
\mathrm d\nu_{s,1}=\mathrm d\rho_s
\]
پیروی می‌کند. در گام آخر، سهم دیفرانسیل
$\nu_{s,H}$ در تابع هدف برابر
$(\mathrm d\nu_{s,H})\alpha_H^{(H)}
\mathsf M_s^{(Y)}\mathbf g_s$ است.
اگر حساسیت کل گام بعد
$\mathbf g_{s,t+1}^{(\nu)}$ باشد، تغییر
$\nu_{s,t}$ علاوه بر سهم مستقیم گام $t$،
از طریق انتقال بعدی سهم
$(\mathrm d\nu_{s,t})P_{\mathcal S_s}
\mathbf g_{s,t+1}^{(\nu)}$ را نیز ایجاد
می‌کند. جمع دو سهم، با استقرا از انتهای افق،
رابطه~\eqref{cmp:eq:ch3-gradient-walk-adjoint}
را نتیجه می‌دهد.
از آنجا که $\nu_{s,t}$ سطری است و حساسیت
ستونی، در ضرب پس‌رونده خود $P_{\mathcal S_s}$
به کار می‌رود، نه ترانهادهٔ آن.

\subsection{ارتباط با بازگشت محاسبهٔ بازنمایی}
\label{app:ch3-gradient-walk-equivalence}

با ضرب بازگشت
رابطه~\eqref{cmp:eq:ch3-matrix-backward-recursion}
از راست در $\mathbf g_s$ داریم:
\[
G_{s,t}\mathbf g_s
=\alpha_t^{(H)}\mathsf M_s^{(Y)}\mathbf g_s
+P_{\mathcal S_s}(G_{s,t+1}\mathbf g_s).
\]
از شرط انتهایی یکسان
$G_{s,H+1}\mathbf g_s=0_{m_s\times1}$،
به استقرای معکوس
رابطه~\eqref{cmp:eq:ch3-gradient-walk-backward-equivalence}
حاصل می‌شود. به‌طور معادل:
\[
\mathbf g_{s,t}^{(\nu)}
=\sum_{u=t}^{H}\alpha_u^{(H)}
P_{\mathcal S_s}^{u-t}\mathsf M_s^{(Y)}\mathbf g_s.
\]
در این محاسبه $\mathbf g_s$ حساسیت ارزیابی‌شده
در نقطهٔ پیش‌رونده است؛ مشتق دوبارهٔ خود آن
برای انتشار مشتق مرتبهٔ اول لازم نیست.

\subsection{مشتق نسبت به توزیع آغاز و ماتریس انتقال}
\label{app:ch3-gradient-walk-transition}

از بازگشت حساسیت داریم:
\[
\begin{aligned}
\alpha_t^{(H)}(\mathrm d\nu_{s,t})
\mathsf M_s^{(Y)}\mathbf g_s
={}&(\mathrm d\nu_{s,t})\mathbf g_{s,t}^{(\nu)}\\
&-(\mathrm d\nu_{s,t})P_{\mathcal S_s}
\mathbf g_{s,t+1}^{(\nu)} .
\end{aligned}
\]
با جایگذاری دیفرانسیل انتقال و جمع روی
$t$، جمله‌های میانی حذف می‌شوند و به دلیل
شرط انتهایی صفر، سهم مبدأ $s$ در دیفرانسیل
تابع هدف برابر است با:
\begin{equation}\label{eq:app-ch3-gradient-walk-loss-differential}
\boxed{
(\mathrm d\mathscr Y_s)\mathbf g_s
=(\mathrm d\rho_s)\mathbf g_{s,1}^{(\nu)}
+\sum_{t=1}^{H-1}\nu_{s,t}
(\mathrm dP_{\mathcal S_s})\mathbf g_{s,t+1}^{(\nu)}.
}
\end{equation}
با استخراج ضریب هر درایهٔ $\mathrm d\rho_s$
و $\mathrm dP_{\mathcal S_s}$،
رابطه~\eqref{cmp:eq:ch3-gradient-walk-transition}
به دست می‌آید. سهم گام‌های متعدد جمع می‌شود
زیرا یک ماتریس انتقال در همهٔ گام‌ها مشترک
است. این مشتق هنوز نسبت به نسخهٔ مختص مبدأ
ماتریس است و اشتراک پارامترها میان مبدأها
در مرحلهٔ بعد لحاظ می‌شود.
صورت ماتریسی مشتق انتقال برابر است با
\[
\nabla_{P_{\mathcal S_s}}\mathcal L_{\mathrm{tr}}
=\sum_{t=1}^{H-1}\nu_{s,t}^{\mathsf T}
(\mathbf g_{s,t+1}^{(\nu)})^{\mathsf T}.
\]
این صورت برای اثبات مفید است، اما در محاسبهٔ
حساسیت سیاست نیازی به تشکیل یا ذخیرهٔ آن نیست.

\subsection{انتقال مشتق به احتمال انتخاب ویژگی}
\label{app:ch3-gradient-walk-local-policy}

از رابطه~\eqref{cmp:eq:ch3-origin-transition-entry}،
با ثابت‌بودن احتمال عمودی و برای یک کاربرد
محلی انتخاب ویژگی داریم:
\[
\mathrm d[P_{\mathcal S_s}]_{(i,a),(j,b)}
=q_b^{(s)}(j\mid i)\,\mathrm d\pi_s(b\mid i,a).
\]
با جایگذاری در دیفرانسیل
رابطه~\eqref{eq:app-ch3-gradient-walk-loss-differential}
و استخراج ضریب $\mathrm d\pi_s(b\mid i,a)$،
حساسیت
\[
g_{s;i,a,b}^{(\pi)}
=\sum_{t=1}^{H-1}\nu_{s,t}(i,a)
\sum_{j:(X_j,b)\in\mathcal S_s}
q_b^{(s)}(j\mid i)\mathbf g_{s,t+1}^{(\nu)}(j,b)
\]
به دست می‌آید. برای نخستین انتخاب، از
رابطه~\eqref{cmp:eq:ch3-block-first-step}
داریم
$\mathrm d\rho_s(j,b)=q_b(j\mid s)
\,\mathrm d\pi_s(b\mid s,\bot)$.
بنابراین ضریب این تغییر در جملهٔ نخست
دیفرانسیل، برابر
\[
g_{s;s,\bot,b}^{(\pi)}
=\sum_{j:(X_j,b)\in\mathcal S_s}
q_b(j\mid s)\mathbf g_{s,1}^{(\nu)}(j,b)
\]
است. بدین ترتیب هر دو رابطهٔ
\eqref{cmp:eq:ch3-gradient-local-policy-middle}
و~\eqref{cmp:eq:ch3-gradient-local-policy-start}
اثبات می‌شوند. این مشتق‌ها مختص کاربرد
محلی احتمال هستند، حتی اگر چند کاربرد از
پارامتر مشترک به دست آمده باشند.

\subsection{گراف تنک، شرایط مرزی و ساختار محاسبات}
\label{app:ch3-gradient-walk-sparse}

با ثابت بودن ساختار گراف تنک و مقصدهای مجاز،
احتمال‌های عمودی تنک‌شده مستقل از پارامترهای
سیاست‌اند. در نتیجه دیفرانسیل انتقال در
گراف تنک نیز همان حاصل‌ضرب احتمال عمودی ثابت
در دیفرانسیل احتمال انتخاب ویژگی است.
برای گام آغاز، احتمال پایهٔ تنک‌شده و برای
گام‌های میانی احتمال اصلاح‌شده پس از حذف
مبدأ به کار می‌رود. با این جایگزینی، تمامی
استدلال‌های پیشین برقرارند. محاسبات مشتق نسبت
به قانون انتقال انتخاب‌شده دقیق‌اند، نه
لزوماً نسبت به گراف کامل پیش از تنک‌سازی.

برای $H=1$، از
$\alpha_1^{(1)}=1$ و شرط انتهایی صفر داریم
$\mathbf g_{s,1}^{(\nu)}
=\mathsf M_s^{(Y)}\mathbf g_s$؛
جمع انتقال‌های میانی تهی است و تنها حساسیت
نخستین انتخاب باقی می‌ماند. در گام‌های
میانی، حساسیت‌های سطرهای بازدیدنشده سهم
صفر دارند، زیرا $\nu_{s,t}(i,a)=0$ است.

برای اجرای ساختاری ضرب پس‌رونده و محاسبهٔ
حساسیت‌ها، در هر گام ابتدا برای هر زوج
$(i,b)$ جمع عمودی
\[
\sum_{X_j\in\mathcal J_{i,b}^{(s)}}
\bar q_b^{(s)}(j\mid i)
\mathbf g_{s,t+1}^{(\nu)}(j,b)
\]
را با زیرروال
\ref{cmp:alg:ch3-vertical-row}
محاسبه می‌کنیم. این جمع مستقل از $a$ است.
ضرب در $\pi_s(b\mid i,a)$ و جمع روی
$b\in A_{X_i}$ بخش انتقالی بازگشت حساسیت
را می‌سازد؛ ضرب همان جمع در
$\nu_{s,t}(i,a)$ سهم گام $t$ در
$g_{s;i,a,b}^{(\pi)}$ را می‌سازد.
بنابراین پیمایش عمودی یک‌بار برای هر
$(i,b)$ و پیمایش افقی روی زوج‌های ویژگی
مجاز کافی است. افزون بر این، طبق ساختار
یک‌داغ ماتریس تجمیع برچسب در
رابطه~\eqref{cmp:eq:ch3-state-aggregation-matrices}،
برای هر حالت $(X_i,a)\in\mathcal S_s$ داریم:
\[
(\mathsf M_s^{(Y)}\mathbf g_s)(i,a)
=\mathbf g_s(y_i).
\]
در پیمایش کامل فضای حالت، محاسبهٔ این بردار
با انتخاب مؤلفهٔ متناظر با برچسب هر حالت،
در $O(m_s)$ عملیات و بدون تشکیل ماتریس
تجمیع ممکن است. اما در اجرای مبتنی بر
تکیه‌گاه، مطابق زیرروال
\ref{cmp:alg:ch3-gradient-walk-backward}،
اصلاً نیازی به ساخت این بردار کامل نیست؛
سهم مستقیم هر حالت فعال $(X_i,a)$ با
خواندن مؤلفهٔ $\mathbf g_s(y_i)$ و ضرب آن
در $\alpha_t^{(H)}$ به دست می‌آید.

با امتیازهای خام متناهی، هر انتخاب مجاز
ویژگی احتمال مثبت دارد. ازاین‌رو اگر
$\nu_{s,t}(i,a)>0$ و انتقال عمودی مجاز
از $(X_i,a)$ به $(X_j,b)$ دارای احتمال
مثبت باشد، آنگاه $\nu_{s,t+1}(j,b)>0$.
بنابراین در بازگشت پس‌رونده، حساسیت هر
حالت فعال تنها به حساسیت مقصدهای فعال
گام بعد وابسته است. با استقرا از گام
$H$ تا گام نخست، حساسیت محاسبه‌شده در
این حالت‌ها دقیقاً با مقدار متناظر در
پیمایش کامل برابر است. سهم حالت‌های
غیرفعال در مشتق‌های موضعی نیز به دلیل
صفر بودن جرم پیش‌رونده صفر است.
در نخستین انتخاب نیز تمام مقصدهای دارای
احتمال مثبت در تکیه‌گاه $\rho_s$ قرار دارند.
پس محدودسازی به تکیه‌گاه‌های دقیق هیچ
تقریبی در مشتق قانون انتقال انتخاب‌شده
ایجاد نمی‌کند.

با فرض دسترسی به سطرهای عمودی و احتمال‌های
سیاست، شمار عملیات هر گام هم‌مرتبه با
محاسبهٔ پیش‌رونده بر همان حالت‌هاست.
برای پیمایش کامل از کران
رابطه~\eqref{cmp:eq:ch3-sparse-linear-bound}
و برای پیمایش مبتنی بر تکیه‌گاه از کران
رابطه~\eqref{cmp:eq:ch3-gradient-reachable-bound}
و تحلیل هزینهٔ زیربخش~
\ref{cmp:subsec:ch3-gradient-procedure}
استفاده می‌کنیم. برای محاسبهٔ مشتق‌های
موضعی می‌توان توزیع‌های پیش‌رونده را
ذخیره یا دوباره محاسبه کرد. حافظهٔ
تاریخچه در پیمایش کامل $O(Hm_s)$ و
در پیمایش مبتنی بر تکیه‌گاه از مرتبهٔ
$O(\sum_{t=1}^{H}|\operatorname{supp}(\nu_{s,t})|)$
است. تنها حساسیت جاری و بعدی نگهداری
می‌شوند؛ در اجرای مبتنی بر تکیه‌گاه،
حافظهٔ این دو متناسب با مجموع اندازهٔ
تکیه‌گاه‌های دو گام متوالی است، جدا از
ساختار گراف، پارامترها و حساسیت‌های
موضعی انباشته.

\section{انتقال مشتق به پارامترهای سیاست}
\label{app:ch3-gradient-policy-parameters}

در این بخش روابط
زیربخش~\ref{cmp:subsec:ch3-gradient-policy-parameters}
را اثبات می‌کنیم. ماسک‌ها ثابت و ناتهی‌اند،
امتیازهای خام مجاز متناهی‌اند و شرایط اعتبار
مشتق‌های مراحل پیشین برقرارند.

\subsection{دیفرانسیل لگاریتم احتمال انتخاب ویژگی}
\label{app:ch3-gradient-log-policy}

برای اختصار، یک سطر موضعی سیاست را با
مجموعهٔ مجاز ثابت $A$، امتیازهای خام
$z=(z_b)_{b\in A}$ و احتمال‌های
$p_b=\exp(z_b)/\sum_{u\in A}\exp(z_u)$
در نظر می‌گیریم. چون $p_b>0$ برای
$b\in A$، داریم:
\[
\log p_b=z_b-\log\sum_{u\in A}\exp(z_u).
\]
با مشتق‌گیری نتیجه می‌شود:
\begin{equation}\label{eq:app-ch3-gradient-log-policy-differential}
\boxed{
\mathrm d\log p_b
=\mathrm dz_b-\sum_{u\in A}p_u\,\mathrm dz_u.
}
\end{equation}
این هویت برای هر ماسک ثابت برقرار است.
در حالت بدون ماسک $A=A_D$ قرار می‌گیرد.
لگاریتم احتمال برای ویژگی نامجاز با
احتمال صفر در این رابطه به کار نمی‌رود.

\subsection{ژاکوبین نرمال‌سازی با بیشینهٔ هموار ماسک‌دار}
\label{app:ch3-gradient-masked-jacobian}

از $\mathrm dp_b=p_b\,\mathrm d\log p_b$
و رابطهٔ پیشین داریم:
\[
\mathrm dp_b=p_b\left(\mathrm dz_b
-\sum_{u\in A}p_u\,\mathrm dz_u\right).
\]
با استخراج ضریب $\mathrm dz_u$ نتیجه می‌شود:
\[
\frac{\partial p_b}{\partial z_u}
=p_b(\mathbf1\{b=u\}-p_u).
\]
با جایگذاری $A=A_{X_i}$ و احتمال‌های
موضعی $\pi_s$، رابطه
\eqref{cmp:eq:ch3-gradient-masked-softmax}
به دست می‌آید. در فضای کامل ویژگی‌ها،
درایه‌های ژاکوبین مربوط به سطر یا ستون
نامجاز صفرند، زیرا احتمال آن ویژگی به
صورت ثابت صفر تعریف می‌شود و امتیازهای
نامجاز نیز در مخرج نرمال‌سازی وارد نمی‌شوند.

\subsection{انتقال مشتق به امتیازهای خام}
\label{app:ch3-gradient-raw-scores}

برای یک تابع اسکالر مشتق‌پذیر از احتمال‌های
همان سطر، با نماد محلی $F(p)$، از قاعدهٔ
زنجیره‌ای و ژاکوبین بالا داریم:
\[
\begin{aligned}
\frac{\partial F}{\partial z_u}
&=\sum_{b\in A}\frac{\partial F}{\partial p_b}
p_b(\mathbf1\{b=u\}-p_u)\\
&=p_u\left[
\frac{\partial F}{\partial p_u}
-\sum_{b\in A}p_b\frac{\partial F}{\partial p_b}
\right].
\end{aligned}
\]
در کاربرد موضعی مورد نظر، حساسیت
$\partial F/\partial p_b$ همان
$g_{s;i,a,b}^{(\pi)}$ است؛ بنابراین
رابطه~\eqref{cmp:eq:ch3-gradient-local-score}
حاصل می‌شود. برای سطر آغاز، $i=s$ و
$a=\bot$ قرار می‌گیرد.
برای محاسبهٔ مشتق تمام امتیازهای یک سطر،
جمع وزن‌دار داخل کروشه یک‌بار محاسبه و
برای همهٔ مؤلفه‌ها بازاستفاده می‌شود.
ازاین‌رو با فرض دسترسی به احتمال‌ها و
حساسیت‌های محلی، هزینهٔ محاسبه $O(|A|)$
است و تشکیل ژاکوبین کامل لازم نیست.

\subsection{خواص گرادیان و جهت‌های بی‌اثر}
\label{app:ch3-gradient-score-invariance}

با جمع کردن مشتق امتیازها روی $u\in A$
و استفاده از $\sum_{u\in A}p_u=1$ داریم:
\[
\begin{aligned}
\sum_{u\in A}\frac{\partial F}{\partial z_u}
&=\sum_{u\in A}p_u\frac{\partial F}{\partial p_u}
-\left(\sum_{u\in A}p_u\right)
\sum_{b\in A}p_b\frac{\partial F}{\partial p_b}\\
&=0.
\end{aligned}
\]
این نتیجه رابطه
\eqref{cmp:eq:ch3-gradient-local-score-zero-sum}
را اثبات می‌کند. افزودن ثابت مشترک به همهٔ
امتیازهای مجاز سطر، احتمال‌های آن را تغییر
نمی‌دهد؛ بنابراین مشتق جهت انتقال مشترک
صفر است. اگر $|A|=1$ باشد، احتمال تنها
ویژگی مجاز برابر یک و تمام مشتق‌های سطر
صفرند. در آرایه‌های امتیاز آزاد، جمع سهم‌های
محلی یک سطر مشترک نیز مجموع مؤلفهٔ صفر
دارد، زیرا مؤلفه‌های نامجاز در هر کاربرد
سهم صفر دارند. این ویژگی برای گرادیان
خام است و لزوماً به تغییر واقعی پارامترها
پس از اعمال بهینه‌سازهایی مانند Adam
منتقل نمی‌شود.

\subsection{تجمیع گرادیان در خانواده‌های امتیاز آزاد}
\label{app:ch3-gradient-free-score-families}

ابتدا همهٔ کاربردهای موضعی امتیازها را
مستقل در نظر می‌گیریم. با جایگذاری
$z_{s;i,a,b}=z_{s;i,a,b}(\theta)$ و
قاعدهٔ زنجیره‌ای، دیفرانسیل تابع هدف
برابر مجموع
$g_{s;i,a,b}^{(z)}\,\mathrm dz_{s;i,a,b}$
روی تمام کاربردهای مجاز است. نخستین
کاربرد در هر گشت از سطر $\bot$ و
نمونهٔ مبدأ، و کاربردهای میانی از حالت‌های
$\mathcal S_s$ به دست می‌آیند. با
جایگذاری دیفرانسیل امتیازها نسبت به
$\theta$، رابطه
\eqref{cmp:eq:ch3-gradient-parameter-chain}
حاصل می‌شود. در امتیازهای آزاد، مشتق
هر کاربرد نسبت به پارامتر مربوط به خود
یک و نسبت به پارامترهای مستقل دیگر صفر
است. در نتیجه کافی است حساسیت کاربردهای
یک پارامتر را جمع کنیم.

در روابط زیر، برای سطرهای میانی
$a,b\in A_D$ و برای سطر آغاز $b\in A_D$
است؛ جمع تهی برابر صفر است. برای
خانواده‌های شرطی به طبقه، قرارداد
$c_s=y_s$ برقرار است.

\paragraph{سیاست ثابت مشترک.}
مطابق رابطه~\eqref{cmp:eq:ch3-fixed-policy-matrix}،
$z_{s;i,a,b}=z_{ab}$ و
$z_{s;s,\bot,b}=z_{\bot b}$ است. بنابراین:
\begin{equation}\label{eq:app-ch3-gradient-shared-scores}
\begin{aligned}
\frac{\partial\mathcal L_{\mathrm{tr}}}{\partial z_{ab}}
&=\sum_{s=1}^n
\sum_{\substack{i\ne s\\a,b\in A_{X_i}}}
g_{s;i,a,b}^{(z)},\\
\frac{\partial\mathcal L_{\mathrm{tr}}}{\partial z_{\bot b}}
&=\sum_{\substack{1\le s\le n\\b\in A_{X_s}}}
g_{s;s,\bot,b}^{(z)}.
\end{aligned}
\end{equation}

\paragraph{سیاست ثابت شرطی به طبقه.}
در این خانواده، امتیاز هر کاربرد از
$Z^{(y_s)}$ انتخاب می‌شود. ازاین‌رو:
\begin{equation}\label{eq:app-ch3-gradient-class-scores}
\begin{aligned}
\frac{\partial\mathcal L_{\mathrm{tr}}}{\partial z^{(c)}_{ab}}
&=\sum_{s\in C_c}
\sum_{\substack{i\ne s\\a,b\in A_{X_i}}}
g_{s;i,a,b}^{(z)},\\
\frac{\partial\mathcal L_{\mathrm{tr}}}{\partial z^{(c)}_{\bot b}}
&=\sum_{\substack{s\in C_c\\b\in A_{X_s}}}
g_{s;s,\bot,b}^{(z)}.
\end{aligned}
\end{equation}

\paragraph{سیاست گره‌محور.}
مطابق رابطه~\eqref{cmp:eq:ch3-policy-family-shapes}،
امتیاز میانی در آرایهٔ $Z^{\mathrm{node}}$
با گره جاری $i$ مشخص می‌شود، نه مبدأ.
برای $a,b\in A_{X_i}$ و به‌ترتیب
$b\in A_{X_i}$ در سطر آغاز داریم:
\begin{equation}\label{eq:app-ch3-gradient-node-scores}
\begin{aligned}
\frac{\partial\mathcal L_{\mathrm{tr}}}
{\partial Z^{\mathrm{node}}_{i,a,b}}
&=\sum_{s\ne i}g_{s;i,a,b}^{(z)},\\
\frac{\partial\mathcal L_{\mathrm{tr}}}
{\partial Z^{\mathrm{node}}_{i,\bot,b}}
&=g_{i;i,\bot,b}^{(z)}.
\end{aligned}
\end{equation}
مشتق درایه‌های نامجاز در هر دو سطر صفر است.

\paragraph{سیاست گره--طبقه.}
برای آرایهٔ $Z^{\mathrm{node,class}}$،
گره جاری و طبقهٔ مبدأ هر دو امتیاز
موضعی را تعیین می‌کنند. برای ویژگی‌های
مجاز همان گره داریم:
\begin{equation}\label{eq:app-ch3-gradient-node-class-scores}
\begin{aligned}
\frac{\partial\mathcal L_{\mathrm{tr}}}
{\partial Z^{\mathrm{node,class}}_{c,i,a,b}}
&=\sum_{\substack{s\in C_c\\s\ne i}}
g_{s;i,a,b}^{(z)},\\
\frac{\partial\mathcal L_{\mathrm{tr}}}
{\partial Z^{\mathrm{node,class}}_{c,i,\bot,b}}
&=\begin{cases}
g_{i;i,\bot,b}^{(z)},&y_i=c,\\
0,&y_i\ne c.
\end{cases}
\end{aligned}
\end{equation}
در آموزش، تنها سطر آغاز مربوط به طبقهٔ
واقعی هر گره به کار می‌رود؛ درایه‌های
آغاز مربوط به دیگر طبقه‌ها و نیز درایه‌های
نامجاز، از تابع هدف فعلی سهم مشتق ندارند.
این امر دربارهٔ سطرهای میانی صدق نمی‌کند؛
زیرا گشت مبدأهای هر طبقه می‌تواند از
گره‌های متعلق به سایر طبقه‌ها عبور کند.

\paragraph{اشتراک مستقل پارامترهای سطر آغاز.}
می‌توان امتیازهای میانی را گره‌محور یا
گره--طبقه نگه داشت، ولی سطر آغاز را به‌صورت
ثابت مشترک یا ثابت شرطی به طبقه تعریف کرد.
در این حالت از رابطهٔ میانی خانوادهٔ
گره‌محور متناظر و از رابطهٔ آغاز
\eqref{eq:app-ch3-gradient-shared-scores}
یا~\eqref{eq:app-ch3-gradient-class-scores}
استفاده می‌شود. شمار جایگاه‌های خام سطر
آغاز مشترک به‌جای $nd$ برابر $d$ و در
حالت شرطی به طبقه به‌جای $Knd$ برابر
$Kd$ است؛ شمار جایگاه‌های میانی تغییر
نمی‌کند. این اعداد شمار درایه‌های خام
آرایه‌اند، نه لزوماً شمار درجات آزادی مؤثر؛
ماسک‌ها و ناوردایی انتقال مشترک نیز بر
درجات آزادی اثر دارند. در نمونهٔ جدید،
همان امتیازهای ثابت آغاز روی ماسک ویژگی‌های
موجود نرمال می‌شوند و درون‌یابی لازم نیست.

\subsection{انتقال مشتق به پارامترهای ماتریس توابع}
\label{app:ch3-gradient-parametric-functions}

برای خانوادهٔ توابع پارامتری در
روابط~\eqref{cmp:eq:ch3-parametric-feature-score}
و~\eqref{cmp:eq:ch3-parametric-start-score}،
با ثابت‌بودن مقادیر ویژگی‌های ورودی، مشتق
امتیاز خام همان مشتق تابع تولیدکنندهٔ
امتیاز نسبت به مؤلفه‌های قابل‌آموزش
$\Theta$ است. برای یک مؤلفهٔ دلخواه
$\theta_j$ از $\theta$ داریم:
\[
\begin{aligned}
\frac{\partial\mathcal L_{\mathrm{tr}}}{\partial\theta_j}
=\sum_{s=1}^{n}\Bigg[
&\sum_{b\in A_{X_s}}g_{s;s,\bot,b}^{(z)}
\frac{\partial F_{c_s,b}
(b(X_s);\Theta)}{\partial\theta_j}\\
&+\sum_{i\ne s}\sum_{a,b\in A_{X_i}}
g_{s;i,a,b}^{(z)}
\frac{\partial F_{c_s,a,b}
(a(X_i),b(X_i);\Theta)}{\partial\theta_j}
\Bigg].
\end{aligned}
\]
اگر مؤلفهٔ $\theta_j$ در چند تابع مشترک
باشد، سهم تمام آن توابع جمع می‌شود؛ اگر
در تابعی حضور نداشته باشد، سهم مشتق آن
تابع صفر است. بنابراین روش به ساختار
داخلی توابع تولیدکنندهٔ امتیاز وابسته نیست
و حساسیت‌های خروجی آن‌ها را می‌توان
به روش مشتق‌گیری مناسب همان ساختار
به پارامترهایشان منتقل کرد. این صورت،
توابع خطی، غیرخطی و شبکه‌های عصبی را
در نقاط دارای مشتق معتبر پوشش می‌دهد.
برای سیاست مستقل از طبقه، اندیس طبقهٔ
مبدأ در انتخاب تابع نقشی ندارد. سیاست
یکنواخت ثابت نیز فاقد پارامتر قابل‌آموزش
است.
